\appendix
\section*{Appendix B: Lean Formalization of Altitude Constraints}

\subsection*{B.0 Overview}

This appendix presents the formal Lean-style mathematical structures
underlying the altitude system of the Constellation–Manifold Toolkit.
The purpose is to define:
\begin{itemize}
    \item altitude identifiers,
    \item allowed and forbidden transitions,
    \item purity constraints,
    \item non-generativity constraints,
    \item boundedness constraints,
    \item violation theorems,
    \item the Series V firewall binding.
\end{itemize}

All definitions are altitude-pure and correspond directly to the
Series IV mathematical semantics and Series V enforcement layer.

\subsection*{B.1 Altitude Type}

We define the altitude type as an inductive set:
\[
\mathcal{A} =
\{A_0, A_1, \Phi_7, D_1, D_2, D_3, \Phi_8, \Phi_9, S_1, I_1\}.
\]

Formally:
\[
\text{Altitude} :=
A_0 \mid A_1 \mid \Phi_7 \mid D_1 \mid D_2 \mid D_3 \mid
\Phi_8 \mid \Phi_9 \mid S_1 \mid I_1.
\]

\subsection*{B.2 Allowed Transition Relation}

Define the allowed transition predicate:
\[
T : \mathcal{A} \times \mathcal{A} \to \{0,1\}.
\]

Allowed transitions:
\[
T(A_0,A_1)=1,\;
T(A_1,\Phi_7)=1,\;
T(\Phi_7,D_1)=1,\;
T(D_1,D_2)=1,\;
T(D_2,D_3)=1,
\]
\[
T(D_3,\Phi_8)=1,\;
T(\Phi_8,\Phi_9)=1,\;
T(\Phi_9,S_1)=1,\;
T(S_1,I_1)=1.
\]

Forbidden transitions:
\[
T(a,b)=0 \quad \text{otherwise}.
\]

\subsection*{B.3 Forbidden Transition Predicate}

Define:
\[
F(a,b) := \neg T(a,b).
\]

This predicate is used by the firewall to detect illegal altitude calls.

\subsection*{B.4 Purity Constraint}

Define the purity predicate:
\[
P : \mathcal{A} \to \{0,1\}.
\]

Purity requires that altitude transformations do not generate new structure:
\[
P(a) = 1 \quad \forall a \in \mathcal{A}.
\]

\subsection*{B.5 Non-Generativity Constraint}

Define:
\[
N : \mathcal{A} \to \{0,1\}.
\]

Non-generativity requires that altitude transformations preserve the domain:
\[
N(a) = 1 \quad \forall a \in \mathcal{A}.
\]

\subsection*{B.6 Boundedness Constraint}

Define:
\[
B : \mathcal{A} \to \{0,1\}.
\]

Boundedness requires:
\[
B(a) = 1 \quad \forall a \in \mathcal{A}.
\]

\subsection*{B.7 Violation Theorem}

For all altitudes $a,b$:
\[
F(a,b) \implies \neg T(a,b).
\]

This is the formal statement that forbidden transitions negate the allowed transition predicate.

\subsection*{B.8 Firewall Binding Theorem}

Define the firewall operator:
\[
\mathrm{Firewall}(a,b,f) =
\begin{cases}
f(), & T(a,b)=1 \land P(b) \land N(b) \land B(b), \\
\mathrm{Violation}, & \text{otherwise}.
\end{cases}
\]

The firewall binding theorem states:
\[
\forall a,b \in \mathcal{A},\;
\mathrm{Firewall}(a,b,f) = f()
\iff
T(a,b)=1 \land P(b) \land N(b) \land B(b).
\]

Otherwise:
\[
\mathrm{Firewall}(a,b,f) = \mathrm{Violation}.
\]

\subsection*{B.9 Altitude Safety Corollary}

If $T(a,b)=0$, then:
\[
\mathrm{Firewall}(a,b,f) = \mathrm{Violation}.
\]

If $P(b)=0$, then:
\[
\mathrm{Firewall}(a,b,f) = \mathrm{Violation}.
\]

If $N(b)=0$, then:
\[
\mathrm{Firewall}(a,b,f) = \mathrm{Violation}.
\]

If $B(b)=0$, then:
\[
\mathrm{Firewall}(a,b,f) = \mathrm{Violation}.
\]

\subsection*{B.10 Collision Invariant Correctness}

The integration altitude $I_1$ produces the collision invariant:
\[
\Omega : X \to \{0,1\}.
\]

Correctness condition:
\[
\Omega(x)=1 \iff \text{system is collision-safe}.
\]

\[
\Omega(x)=0 \iff \text{violation detected}.
\]

This completes the formal Lean-style appendix.
